\begin{center}\begin{tikzpicture}[font=\small,>=stealth,line join=round]
\draw[gray!40,rounded corners=3pt] (0,0) rectangle (14.2,8.3);\node at (7.1,7.700000000000001) {从指数反求次数：真数只能为正，次数却没有正负限制};
\begin{scope}[shift={(1.6660377358490566,4.1269230769230765)},x=1.1320754716981132cm,y=0.7076923076923076cm]\draw[->,gray] (-0.5,0)--(4.8,0) node[right]{$x$};\draw[->,gray] (0,-3.5)--(0,3) node[above]{$y$};\draw[blue!75!black,thick,solid,domain=0.125:4.2,samples=180] plot(\x,{ln(\x)/ln(2)});\fill[orange!85!black] (0.25,-2) circle(2.2pt);\fill[orange!85!black] (1,0) circle(2.2pt);\fill[orange!85!black] (4,2) circle(2.2pt);\draw[gray,dashed] (.05,-2)--(4.3,-2);\draw[gray,dashed] (.05,2)--(4.3,2);\draw[gray!35,densely dotted] (1,0)--(1,-3.5);\node[below=3pt] at (1,-3.5){$1$};\draw[gray!35,densely dotted] (4,0)--(4,-3.5);\node[below=3pt] at (4,-3.5){$4$};\end{scope}\node[anchor=west,align=left,text width=5.0cm] at (8.4,4.0) {$x=2^y$总大于0\\[8pt]$y=-2\Rightarrow x=\frac14$\\[8pt]$y=0\Rightarrow x=1$\\[8pt]$y=2\Rightarrow x=4$\\[8pt]任意实数$y$都有正数$x$\\[8pt]定义域$(0,+\infty)$，值域$\mathbb R$};
\node[text width=12.6cm,align=center] at (7.1,.65) {水平线y=c与曲线的交点横坐标为2的c次方，总能找到且总为正数。};
\end{tikzpicture}\end{center}
